% Level 2, Week 11 homework -- project work
\homeworkweek{11}{Project Studio 2}{15--20 minutes. Do Core first; Challenge and Creative are optional.}
\namedateline

\begin{core}[Three passing tests and one fix (about 10 minutes)]
\begin{enumerate}
  \item Open your project folder in Thonny. Make sure you have a file called \texttt{test\_mygame.py} (or \texttt{test\_} plus your game's name) with \textbf{at least 3 tests}.
  \item Run your game once with \keys{F5}, close its window, then \menu{Tools > Open system shell…} and type \texttt{python -m pytest -v} (or \texttt{python3 -m pytest -v}).
  \item Copy the last line of the output here (for example \enquote{3 passed in 0.05s}): \answerlines{1}
  \item Pick the top MUST line from your playtest bug list and fix it. What did you fix? \answerlines{1}
\end{enumerate}
\end{core}

\begin{challenge}[Test first]
Find a bug or a missing feature in a logic function. Write the test \textbf{first} and run it: it
should say FAILED. Then change your code until it says PASSED. Write your test here:
\answerspace{3cm}
Why is it useful to see the test fail first? \answerlines{1}
\end{challenge}

\begin{creative}[How to play]
Add a \enquote{how to play} message to your project's screen (with \texttt{write} and a font size of
18 or more), or write it on a card to place next to the computer at the showcase. Include: the goal,
the keys, and how the game ends (or, for an Art Studio, how to draw, save and load). Test it on a
family member: can they start without asking you anything?
\answerlines{3}
\end{creative}

\begin{selfcheck}
  \item At least 3 of my tests say PASSED.
  \item My project saves something to a file and loads it again.
  \item I fixed at least one bug from the playtest.
  \item I ticked the accessibility checklist in workbook chapter 11 (colours, shapes, text size, keys explained, no flashing).
  \item High-score names and saved files use nicknames only.
\end{selfcheck}
